\section{Linear thermoelastic plate model}

Consider a homogeneous isotropic plate occupying $-h\le z\le h$, unbounded in
$x$, with displacement $\bm u=(u_x,u_z)$. In the small-strain approximation,
\begin{align}
  \epsilon_{ij} &= \frac{1}{2}(\partial_i u_j+\partial_j u_i),\\
  \sigma_{ij} &= \lambda\epsilon_{kk}\delta_{ij}
  +2\mu\epsilon_{ij}-(3\lambda+2\mu)\alpha\Delta T\,\delta_{ij},\\
  \rho\,\ddot u_i &= \partial_j\sigma_{ij},
\end{align}
where $\lambda$ and $\mu$ are the Lam\'e constants, $\rho$ is density, and
$\alpha$ is the coefficient of linear thermal expansion. The free surfaces obey
\begin{equation}
  \sigma_{xz}(x,\pm h,t)=\sigma_{zz}(x,\pm h,t)=0.
\end{equation}
The bulk longitudinal and transverse speeds are
\begin{equation}
  c_L=\sqrt{\frac{\lambda+2\mu}{\rho}},\qquad
  c_T=\sqrt{\frac{\mu}{\rho}}.
\end{equation}

The temperature rise is governed, at the simplest level, by
\begin{equation}
  \rho C_p\,\partial_t\Delta T-K\nabla^2\Delta T=Q(x,z,t),
  \label{eq:heat}
\end{equation}
with volumetric heat capacity $\rho C_p$, conductivity $K$, and absorbed optical
power density $Q$. Equations~\eqref{eq:heat} and the elastodynamic equations are
one-way coupled when mechanical heating is negligible.
